\textbf{COLUNI 2026 -- Questão 26}\par\medskip O Brasil sediará, em novembro de 2025, a 30ª Conferência das Nações Unidas sobre Mudança do Clima (COP30), em Belém (PA), reunindo representantes de mais de 190 países para avançar nas negociações sobre a agenda climática global. Desde a assinatura do Acordo de Paris (2015), países do mundo todo têm buscado reduzir suas emissões de $CO_2$ para conter o aquecimento global, tema que ganhará ainda mais destaque com a realização da COP 30 no Brasil.\par
Disponível em: https://www.gov.br/trabalho-e-emprego/pt-br/assuntos/cop-30-brasil-amazonia#:~:text=O%20Brasil%20sediar% C3% A1%2C%20em%20novembro,sobre%20a%20agenda%20clim% C3% A1tica%20global. Acesso em: 26.jul.2025.

O $CO_2$ é um produto da combustão de combustíveis fósseis, como a gasolina, que é um líquido menos denso que a água e, quando misturada a ela, forma duas fases distintas. Um de seus principais componentes, o octano ($C_8H_{18}$), sofre combustão completa segundo a reação:
\[
2\,C_8H_{18}(l)+9\,O_2(g)\longrightarrow16\,CO_2(g)+18\,H_2O(l)
\]
Com base nas informações fornecidas e nos conhecimentos sobre ciências, é CORRETO afirmar que:\begin{enumerate}
\item a gasolina afunda na água porque sua densidade é maior que 1,0 g/mL.
\item a combustão do octano é uma reação que libera calor.
\item o octano reage com o gás oxigênio por neutralização formando $CO_2$.
\item a filtração simples é utilizada para separar gasolina e água.\end{enumerate}